\section{Background}
\label{sec:background}

\para{Legal foundations and user expectations of data deletion}
%
Data deletion is an important mechanism for protecting user privacy, supported by both legal requirements and user expectations. 
%
For example, Article 17 of the GDPR establishes the right to erasure under specified conditions, commonly known as the right to be forgotten~\cite{gdpr2016}. 
%
Similarly, the CCPA gives California residents the right to request deletion of their personal information~\cite{ccpa2018}. 
%
Beyond legal requirements, users expect to have control over their personal information, including the ability to delete previously shared data~\cite{murillo2018if}.
%
As LLM agents accumulate substantial amounts of personal information through ongoing conversations (e.g., voluntary user disclosures and requests for user information to personalize their responses) and may also use user data to train their models, understanding whether and how data deletion works in these systems becomes critical to protecting users' privacy rights.

\para{Data deletion methods of LLM agents}
%
Looking through user-facing features, LLM agents may retain user-provided information across different locations, such as conversation histories, saved memories\footnote{Here, saved memory refers to a user-facing feature that retains information across conversations, rather than information encoded in model parameters or held in the context window.}, and account-level profiles or settings.
%
Accordingly, LLM agents offer different data deletion methods that appear to target specific data storage locations, while their effects on data retained elsewhere are not always clear.
%
For example, deleting a conversation removes it from the user-visible conversation history, but its effects on information retained in saved memories or used for model training remain unclear.


To understand the deletion methods available to users, we conducted a systematic walkthrough of six widely used commercial LLM agents using free consumer accounts:  ChatGPT, Claude, Gemini, Copilot, Perplexity, and DeepSeek.
%
Following prior research on LLM privacy controls~\cite{li2026privacy}, we inspected each agent's interface and settings to identify the available data deletion methods.
%
As shown in Table~\ref{tab:methods}, we identified six data deletion methods: (D1) sending a deletion-request prompt through conversations; (D2) deleting a single conversation or (D3) clearing all conversation history through UI buttons; (D4) deleting a specific saved memory or (D5) clearing all saved memories through memory settings; and (D6) using an account-level privacy dashboard.
%
These methods differ in their intended scope, ranging from individual conversations and saved facts to broader conversation histories and account-level data management.
%
Their availability also varies across agents, as detailed in Table~\ref{tab:methods}.


\begin{table}[t]
\centering
\caption{Availability of six data deletion methods across commercial LLM agents. A dash indicates that a method was unavailable on the tested free account.}
\label{tab:methods}
\footnotesize
\setlength{\tabcolsep}{3pt}
\begin{tabular}{@{}p{0.5\columnwidth}cccccc@{}}
\toprule
Data deletion method
& \rotatebox{75}{ChatGPT}
& \rotatebox{75}{Claude}
& \rotatebox{75}{Gemini}
& \rotatebox{75}{Copilot}
& \rotatebox{75}{Perplexity}
& \rotatebox{75}{DeepSeek} \\
\midrule
(D1) Send a deletion-request prompt
& \ding{51} & \ding{51} & \ding{51} & \ding{51} & \ding{51} & \ding{51} \\
(D2) Delete a single conversation
& \ding{51} & \ding{51} & \ding{51} & \ding{51} & \ding{51} & \ding{51} \\
(D3) Clear all conversation history
& \ding{51} & -- & \ding{51} & -- & \ding{51} & \ding{51} \\
(D4) Delete a specific saved memory
& -- & \ding{51} & \ding{51} & \ding{51} & -- & -- \\
(D5) Clear all saved memories
& \ding{51} & \ding{51} & \ding{51} & \ding{51} & -- & -- \\
(D6) Use a privacy dashboard
& -- & -- & -- & \ding{51} & -- & -- \\
\bottomrule
\end{tabular}
\end{table}